\begin{lemma}[Stability of the rescaled equation as its parameter degenerates]
Let $F$ be a finite-dimensional real inner-product space, let $\Omega$ be an
open subset of the space-time $\ell^2(F\times\R)$, and let
$(\lambda_j)_{j\ge0}$ be real numbers with $\lambda_j\to0$. Let
$(v_j)_{j\ge0}$ be real functions on the space-time such that each $v_j$ is a
viscosity solution of $D^2v_j[(\nabla_xv_j,1+\lambda_j\partial_tv_j)]=0$ on
$\Omega$ in the sense of \noderef{IsRescaledSolutionOn}, and let $v_*$ be a real
function on the space-time such that $v_j\to v_*$ uniformly on every compact
subset of $\Omega$. Then $v_*$ is a viscosity solution of $F[v_*]=0$ on
$\Omega$ in the sense of \noderef{IsLimitSolutionOn}.
\end{lemma}
\begin{proof}
Write $R_\lambda[\phi](\zeta)$ for the rescaled operator
$D^2\phi(\zeta)[(\nabla_x\phi(\zeta),1+\lambda\partial_t\phi(\zeta))]$ of
\noderef{RescaledOp} and $F[\phi](\zeta)$ for the limiting operator
$D^2\phi(\zeta)[(\nabla_x\phi(\zeta),1)]$ of \noderef{LimitOp}. Note the two
agree at $\lambda=0$:
\begin{equation}\label{eq:rls-zero}
  R_0[\phi](\zeta)=F[\phi](\zeta),
\end{equation}
because $1+0\cdot\partial_t\phi(\zeta)=1$ and the two operators are then the
second derivative of $\phi$ at $\zeta$ evaluated on the same pair of vectors.

\emph{Step 0: continuity of the limit.}
Let $\zeta\in\Omega$. Since $\Omega$ is open there is $\varrho>0$ with
$\overline B_\varrho(\zeta)\subseteq\Omega$, a compact set, so $v_j\to v_*$
uniformly on it. By \noderef{IsRescaledSolutionOn} each $v_j$ is continuous on
$\Omega$, hence on $\overline B_\varrho(\zeta)$, and a uniform limit of
continuous functions on a set is continuous on that set. Thus $v_*$ is
continuous at every point of $\Omega$, i.e.\ continuous on $\Omega$.

\emph{Step 1: strictifying an upper test.}
Let $\zeta_0\in\Omega$ and let $\phi$ be $C^2$ near $\zeta_0$ with
$\zeta\mapsto v_*(\zeta)-\phi(\zeta)$ having a local maximum on $\Omega$ at
$\zeta_0$. Set
\[
  \phi^\sharp(\zeta)=\phi(\zeta)+|\zeta-\zeta_0|^4 .
\]
The function $\zeta\mapsto|\zeta-\zeta_0|^4=\bigl(|\zeta-\zeta_0|^2\bigr)^2$ is a
polynomial in the coordinates, hence smooth, and
\begin{equation}\label{eq:rls-fourth}
  \nabla\bigl(|\cdot-\zeta_0|^4\bigr)(\zeta_0)=0,
  \qquad
  D^2\bigl(|\cdot-\zeta_0|^4\bigr)(\zeta_0)=0 ,
\end{equation}
since its gradient is $4|\zeta-\zeta_0|^2(\zeta-\zeta_0)$, which vanishes at
$\zeta_0$ together with its derivative. Hence $\phi^\sharp$ is $C^2$ near
$\zeta_0$ with $\nabla\phi^\sharp(\zeta_0)=\nabla\phi(\zeta_0)$ and
$D^2\phi^\sharp(\zeta_0)=D^2\phi(\zeta_0)$, so by \noderef{LimitOp}
\begin{equation}\label{eq:rls-same-jet}
  F[\phi^\sharp](\zeta_0)=F[\phi](\zeta_0).
\end{equation}

Choose $\varrho>0$ so small that $\overline B_\varrho(\zeta_0)\subseteq\Omega$,
that $\phi$ is $C^2$ on an open set $V$ with
$\overline B_\varrho(\zeta_0)\subseteq V\subseteq\Omega$, and that $v_*-\phi$
attains its maximum over $\overline B_\varrho(\zeta_0)$ at $\zeta_0$; the last is
possible because the maximum is local on $\Omega$ and $\Omega$ is open. Then for
$\zeta\in\overline B_\varrho(\zeta_0)$ with $\zeta\ne\zeta_0$,
\[
  v_*(\zeta)-\phi^\sharp(\zeta)
   =\bigl(v_*(\zeta)-\phi(\zeta)\bigr)-|\zeta-\zeta_0|^4
   <v_*(\zeta_0)-\phi(\zeta_0)=v_*(\zeta_0)-\phi^\sharp(\zeta_0),
\]
so $\zeta_0$ is the \emph{unique} maximum point of $v_*-\phi^\sharp$ on
$\overline B_\varrho(\zeta_0)$.

\emph{Step 2: nearby maxima of the approximations.}
The set $\overline B_\varrho(\zeta_0)$ is compact and contained in $\Omega$, and
each $v_j-\phi^\sharp$ is continuous on it, so it attains a maximum at some
$\zeta_j\in\overline B_\varrho(\zeta_0)$. We claim $\zeta_j\to\zeta_0$. By
compactness it suffices to show that every convergent subsequence has limit
$\zeta_0$. Let $\zeta_{j_k}\to\bar\zeta\in\overline B_\varrho(\zeta_0)$. Uniform
convergence of $v_j$ to $v_*$ on $\overline B_\varrho(\zeta_0)$ and continuity of
$v_*-\phi^\sharp$ there give
\[
  (v_{j_k}-\phi^\sharp)(\zeta_{j_k})\longrightarrow
  (v_*-\phi^\sharp)(\bar\zeta).
\]
On the other hand, for every $\zeta\in\overline B_\varrho(\zeta_0)$,
$(v_{j_k}-\phi^\sharp)(\zeta_{j_k})\ge(v_{j_k}-\phi^\sharp)(\zeta)
 \to(v_*-\phi^\sharp)(\zeta)$, so
$(v_*-\phi^\sharp)(\bar\zeta)\ge(v_*-\phi^\sharp)(\zeta)$ for all such $\zeta$.
Thus $\bar\zeta$ is a maximum point of $v_*-\phi^\sharp$ on
$\overline B_\varrho(\zeta_0)$, and by the uniqueness in Step 1,
$\bar\zeta=\zeta_0$. Hence $\zeta_j\to\zeta_0$; in particular
$|\zeta_j-\zeta_0|<\varrho$ for all large $j$, so $\zeta_j$ is an interior point
of $\overline B_\varrho(\zeta_0)$ and therefore $v_j-\phi^\sharp$ has a local
maximum on $\Omega$ at $\zeta_j$.

\emph{Step 3: passing the inequality, in the point and in the parameter.}
For such $j$, $\phi^\sharp$ is $C^2$ near $\zeta_j$ and $v_j-\phi^\sharp$ has a
local maximum on $\Omega$ at $\zeta_j$, so the subsolution clause of
\noderef{IsRescaledSolutionOn} for $v_j$ gives
\begin{equation}\label{eq:rls-at-j}
  R_{\lambda_j}[\phi^\sharp](\zeta_j)\ge0 .
\end{equation}
It remains to pass to the limit in \emph{both} arguments, the point $\zeta_j$ and
the parameter $\lambda_j$. Since $\phi^\sharp$ is $C^2$ on the open set $V$, both
$\zeta\mapsto\nabla\phi^\sharp(\zeta)$ and
$\zeta\mapsto D^2\phi^\sharp(\zeta)$ are continuous on $V$; write
$\nabla_x\phi^\sharp(\zeta)$ and $\partial_t\phi^\sharp(\zeta)$ for the two
components of $\nabla\phi^\sharp(\zeta)$, which are therefore continuous on $V$
as well. For large $j$ we have $\zeta_j\in\overline B_\varrho(\zeta_0)\subseteq V$,
and $\zeta_j\to\zeta_0$, so
\[
  \nabla_x\phi^\sharp(\zeta_j)\to\nabla_x\phi^\sharp(\zeta_0),\qquad
  \partial_t\phi^\sharp(\zeta_j)\to\partial_t\phi^\sharp(\zeta_0),\qquad
  D^2\phi^\sharp(\zeta_j)\to D^2\phi^\sharp(\zeta_0),
\]
the last in the space of continuous symmetric bilinear forms. Since
$\lambda_j\to0$ and $\partial_t\phi^\sharp(\zeta_j)$ converges, hence is bounded,
\[
  1+\lambda_j\partial_t\phi^\sharp(\zeta_j)\longrightarrow1,
\]
and therefore the vectors
$\xi_j=\bigl(\nabla_x\phi^\sharp(\zeta_j),1+\lambda_j
  \partial_t\phi^\sharp(\zeta_j)\bigr)$
converge to $\xi_0=\bigl(\nabla_x\phi^\sharp(\zeta_0),1\bigr)$. Evaluation of a
continuous bilinear form on a pair of vectors is jointly continuous in the form
and the vectors, so by \noderef{RescaledOp} and \noderef{LimitOp}
\[
  R_{\lambda_j}[\phi^\sharp](\zeta_j)
   =D^2\phi^\sharp(\zeta_j)[\xi_j]
   \longrightarrow D^2\phi^\sharp(\zeta_0)[\xi_0]
   =F[\phi^\sharp](\zeta_0).
\]
Passing to the limit in \eqref{eq:rls-at-j} gives
$F[\phi^\sharp](\zeta_0)\ge0$, and \eqref{eq:rls-same-jet} yields
$F[\phi](\zeta_0)\ge0$, which is the subsolution inequality of
\noderef{IsLimitSubOn} at $\zeta_0$.

\emph{Step 4: lower tests.}
If instead $v_*-\phi$ has a local minimum on $\Omega$ at $\zeta_0$, repeat
Steps 1--3 with $\phi^\sharp(\zeta)=\phi(\zeta)-|\zeta-\zeta_0|^4$, with minima
in place of maxima, and with the supersolution clause of
\noderef{IsRescaledSolutionOn}, which gives
$R_{\lambda_j}[\phi^\sharp](\zeta_j)\le0$ in place of \eqref{eq:rls-at-j}. The
limit passage of Step 3 is unchanged, identity \eqref{eq:rls-fourth} again gives
$F[\phi^\sharp](\zeta_0)=F[\phi](\zeta_0)$, and the conclusion is
$F[\phi](\zeta_0)\le0$, the supersolution inequality of
\noderef{IsLimitSuperOn} at $\zeta_0$.

Together with Step 0 and the openness of $\Omega$, this verifies both clauses of
\noderef{IsLimitSolutionOn} for $v_*$.
\end{proof}
